\chapter{Coordinate Frame and Rotation}\label{ch:frame}

A global ocean mesh must contend with the convergence of the meridians at the geographic poles, where
the longitude metric $\zzero=\cos\latx\,\Mlat$ of \cref{sec:intro-conv} vanishes and longitude becomes
ill-conditioned. FESOM2 meshes commonly avoid this by rotating the computational frame so that its
poles lie over land --- conventionally over Greenland and Antarctica --- leaving the entire ocean far from
the coordinate singularity. This chapter describes how the engine supports such a rotated frame: the
rotation itself, the distinction between rotating points and rotating vectors, what is and is not
affected, the contract that keeps the rotation invisible to the user, and the accuracy it achieves. The
facility is available on FESOM2 and HEALPix meshes; FESOM-C meshes and NEMO-derived meshes reject a
non-null \conf{mesh.rotation} with a configuration error.

\section{The rotated frame}\label{sec:frame-rotation}

The rotation is configured by the key \conf{mesh.rotation}. Its default is \conf{null}, a geographic
frame in which no rotation is applied and the trajectories are computed directly in longitude and
latitude; a run with this default is byte-identical to one on an engine without the facility. Setting
\conf{mesh.rotation} to \conf{fesom2} selects FESOM2's own Euler angles $(\alpha,\beta,\gamma)=
(50^\circ,15^\circ,-90^\circ)$, which place the computational poles in Greenland and in East Antarctica;
an explicit triple of angles may also be given. The rotation is a rigid map of the sphere,
\begin{equation}\label{eq:frame-map}
  \mathbf{r}' = \mathbf{R}(\alpha,\beta,\gamma)\,\mathbf{r},\qquad
  \mathbf{r}=(\cos\latx\cos\lonx,\ \cos\latx\sin\lonx,\ \sin\latx),
\end{equation}
with $\mathbf{R}$ the Euler rotation matrix and $\mathbf{R}^{\mathsf T}$ its inverse; rotated longitudes are
wrapped to $[-180^\circ,180^\circ)$. It is applied once when the mesh and each velocity snapshot are read,
never inside the per-particle loop. With the \conf{fesom2} angles the rotated FESOM2 mesh reproduces the
model's own element areas to machine precision.

\section{Rotating points and rotating vectors}\label{sec:frame-vectors}

The rotation acts differently on positions and on velocities, and the distinction is essential to
correctness (\cref{fig:rotation}). Node coordinates and element centres are \emph{points}: they are
mapped by \cref{eq:frame-map}. The element centres are rotated as points and not recomputed from the
rotated nodes, because a centroid computed on the sphere does not commute with the rotation. The
horizontal velocity components, by contrast, are the components of a \emph{tangent vector} expressed in
the local east--north frame, and that local frame itself turns from place to place; the velocity must
therefore be re-expressed in the rotated frame's local east--north basis. With $c$ and $s$ the cosine
and sine of the local angle through which the unit-east direction turns, the transformed components are
\begin{equation}\label{eq:vec-rot}
  u' = c\,u - s\,v, \qquad v' = s\,u + c\,v,
\end{equation}
a rotation that preserves the speed, $u'^2+v'^2=u^2+v^2$. The factors $(c,s)$ are the rotated-frame
components of $\mathbf{R}\,\hat{\mathbf{e}}$, with $\hat{\mathbf{e}}=(-\sin\lonx,\cos\lonx,0)$ the geographic
unit east vector, computed once per velocity carrier --- per node, per element centre, or both, as the
mesh's velocity layout requires --- from the \emph{geographic} coordinates, before the points are
rotated. At an exact pole the east direction is taken as the limit at the stored longitude.

Scalar and vertical quantities have no horizontal direction and pass through unchanged: the vertical
velocity, the sea-surface height, the diffusivity and the sampled scalars. For the diffusivity one
qualification applies: a diffusivity whose zonal and meridional values differ acts along the rotated
axes rather than along east and north, and the engine warns once when it meets such a pair. The
stochastic kicks move along great circles (\cref{ch:diffusion}), so an isotropic diffusivity gives the
same statistics in either frame.

\begin{figure}[t]
\centering
\begin{tikzpicture}[scale=1.3]
  \coordinate (O) at (0,0);
  \def\th{28}
  \draw[-{Stealth[length=5pt]},gray!70,line width=0.8pt] (O)--(2.4,0) node[right,font=\footnotesize,black]{east};
  \draw[-{Stealth[length=5pt]},gray!70,line width=0.8pt] (O)--(0,2.4) node[above,font=\footnotesize,black]{north};
  \draw[-{Stealth[length=5pt]},lagelem,line width=0.9pt,densely dashed] (O)--(\th:2.4) node[right,font=\footnotesize,lagelem!70!black]{east$'$};
  \draw[-{Stealth[length=5pt]},lagelem,line width=0.9pt,densely dashed] (O)--({\th+90}:2.4) node[above,font=\footnotesize,lagelem!70!black]{north$'$};
  \draw[->,black] (1.5,0) arc (0:\th:1.5);
  \node[font=\footnotesize] at ({\th/2}:1.72) {$\theta$};
  \draw[-{Stealth[length=6pt]},lagblue,line width=1.3pt] (O)--(62:2.0) node[above right,font=\footnotesize,lagblue!70!black]{$(u,v)$};
  \node[font=\footnotesize] at (1.0,-0.5) {$u'=c\,u-s\,v,\quad v'=s\,u+c\,v\quad(c=\cos\theta,\ s=\sin\theta)$};
\end{tikzpicture}
\caption{Rotation of a horizontal velocity into the rotated frame. The velocity $(u,v)$ is a tangent
vector; its components are re-expressed in the rotated local east--north basis (dashed), turned through
the local angle $\theta$, by \cref{eq:vec-rot}. The map preserves the speed. Node coordinates and element
centres, being points rather than vectors, are instead mapped directly by the rotation.}
\label{fig:rotation}
\end{figure}

\section{The user-facing contract}\label{sec:frame-contract}

The rotation is arranged to be invisible outside the engine. Release positions are always supplied in
geographic coordinates and are rotated into the computational frame when the run starts. Trajectory
output is always geographic: at every record the positions and the horizontal velocity diagnostics are
un-rotated on a copy of the state, while the internal state stays in the rotated frame; the vertical
velocity and the sampled scalars are written as they are. The release positions written with the
trajectories (\cref{sec:out-schema}) are geographic too: the particle state carries them in the rotated
frame, checkpoints store them in that frame, and they are un-rotated by the same rule when they are
written. A restart checkpoint records the three angles
of the frame in which it was written, and a resuming driver validates them against the mesh
(\fn{\_validate\_restart\_frame}, \cref{ch:output}), so a run is not continued in a different frame by
accident; the output file carries the angles as metadata whenever
the frame is rotated. Within this contract the rotation is purely an internal device for conditioning
the arithmetic near the poles. Periodic fold windows (\cref{ch:lifecycle}) are defined in the mesh's
coordinates and cannot be combined with a rotation; the combination is refused.

\section{Rotation on HEALPix meshes}\label{sec:frame-healpix}

A FESOM2 model builds its mesh in its own rotated frame; a HEALPix product carries geographic east and
north components on a global pixel grid. Rotation of a HEALPix mesh is therefore a pure runtime
transform: node coordinates are rotated as points, the velocities receive tangent factors per node (per pixel centre on the pixel mesh), the
search structures are rebuilt on the rotated coordinates, and the tracking layer handles the rest as for
FESOM2. The rotated frame has two poles of its own, and the polar pathologies move there, so the loader
guards against them: a wet element that straddles a rotated pole is refused with an error that names
where the poles land, and wet nodes poleward of $88^\circ$ rotated latitude produce a warning. With the
\conf{fesom2} angles both poles fall on land, and on the global HEALPix products tested the wet nodes stay
below about $85^\circ$ rotated latitude. In the geographic frame, by contrast, the containment test rejects
points poleward of the first ring of pixel centres, $\arcsin\bigl(1-1/(3\,n_{\text{side}}^2)\bigr)$ ---
$84.15^\circ$ at $n_{\text{side}}=8$, $89.91^\circ$ at $n_{\text{side}}=512$ --- so a particle there dies
as \code{OUTSIDE}; for global HEALPix runs the rotation is the recommended setting.

\section{Accuracy and when rotation is needed}\label{sec:frame-accuracy}

Rotation is a change of chart, not of physics, and away from the poles a correctly implemented run
differs between the geographic and rotated frames only at the level of the discretization. In the
engine's verification the two frames agree to about $10^{-6}$ degrees at mid-latitudes, and the
difference shrinks as the step is refined: at second order in the step for element-layout velocity and
at first order for node-layout velocity. Across a pole only the rotated frame can track: a meridional
flow over the pole follows the analytic great circle to within \SI{7e-8}{\metre}, whereas the geographic
frame breaks down there. A genuine framing error would produce discrepancies of tens of degrees,
independent of the step. Over long integrations the two frames diverge like any two slightly different
discretizations of a chaotic flow; their statistics, not their individual paths, are what agree.

Rotation becomes necessary only when the mesh carries wet points close to a geographic pole. The
FESOM2 reader warns, and suggests the \conf{fesom2} setting, when a mesh extends beyond about
$\pm 88^{\circ}$ without a rotation, and warns about a node exactly at a pole when rotated; the FESOM-C
reader warns without the suggestion; the HEALPix reader warns as described in \cref{sec:frame-healpix}.
For regional domains away from the poles the geographic frame is appropriate and is the default.

\codeanchor{The maps are \fn{euler\_matrix}, \fn{rotate\_lonlat}/\fn{unrotate\_lonlat} and
\fn{rotate\_vectors}/\fn{unrotate\_vectors} in \srcfile{core/rotation.py}; the angles and vector factors
are \fn{rotation\_angles} and \fn{\_vec\_rot\_factors} (\srcfile{core/mesh.py}); the readers apply them
in \fn{\_apply\_mesh\_rotation} (\srcfile{readers/fesom2.py}) and \fn{\_apply\_rotation}
(\srcfile{readers/healpix.py}); the tracking layer rotates seeds in \fn{\_maybe\_rotate\_seeds},
un-rotates records in \fn{\_geo\_view\_for\_record} and checks restarts in
\fn{\_validate\_restart\_frame} (\srcfile{tracking.py}).}
